
\needspace{4\baselineskip}
\color{uniblue}{\subsection{Защита от потери охлаждения (ЗПО)}\label{sec:zpo}}
\color{black}

\begin{enumerate}

\item
ЗПО служит для формирования сигнала на отключение трансформатора при отказе системы охлаждения трансформатора, дополнительными контролируемыми факторами для формирования сигнала отключения могут являться сигналы срабатывания токовых органов ЗПО (см. подраздел \ref{sec:tozpo}) и датчика температуры масла. В устройстве защита от потери охлаждения представлена функцией «ЗПО».

\item
Трансформатор с высшим напряжением 35 кВ может оснащаться системой охлаждения типа «Д» (естественная циркуляция масла с принудительной циркуляцией воздуха). Для таких трансформаторов необходимо предусматривать автоматическое управление пуском и остановом вентиляторов обдува. В случае неисправности системы охлаждения и превышения нормальных параметров эксплуатации трансформатора (повышение температуры верхних слоев масла, превышение номинального тока трансформатора) в некоторых случаях может потребоваться отключение трансформатора. 

\item
Ввод функции в работу на этапе параметрирования устройства осуществляется программным переключателем SGF1 «Ввод\_функции». Информация о введенной функции отображается выходным сигналом «ЗПО: Ввод». 

\item
Функция может быть выведена из работы оперативно путем активации сигнала «ОВ ЗПО», что характеризуется активным состоянием сигнала «ЗПО: Оперативный вывод» на выходе функции. 

\item
При активном состоянии входного сигнала «ЗПО~на~сигн» действие отключающей логики функции переводится на сигнал.

\item
Пуск функции происходит при приеме сигнала «Отказ сист. охл.» от системы охлаждения, формирование сигнала отключения происходит через выдержку времени таймера T1 (уставка «Тср»). При этом возможен контроль выполнения следующих условий:
\begin{itemize}
\item срабатывание токовых органов по сторонам ВН, НН (возможная уставка срабатывания -- $\textsl{1,05}I_\textsl{ном}$), управляется программным переключателем SGF2 «Контр\_тока»;
\item срабатывание датчика температуры (возможная уставка срабатывания -- $\textsl{55}^\circ C$), управляется программным переключателем SGF3 «Контр\_темп\_масла».
\end{itemize}
\par {\small{ \textls[100]{Примечание} -- Возможные уставки срабатывания приведены в качестве примера из ГОСТ Р 52719-2007, при настройке функции необходимо руководствоваться, в том числе и требованиями завода-изготовителя трансформатора.}}

\item
Алгоритм работы функции выполнен в соответствии с рисунком \ref{zpo:img1}. В таблице \ref{app_set:zpo} приведены параметры, необходимые для настройки функции.

\vspace{1mm}
\begin{figure}[H]
\centering
\includegraphics[width=0.8\textwidth,height=0.8\textheight,keepaspectratio]{\fbpath/08. ЗПО Т 35/_latex/img/1.pdf}
\captionof{figure}{Функциональная схема логики ЗПО}\label{zpo:img1}
\end{figure}

\end{enumerate}
\FloatBarrier % Все float'ы ДО этой строки будут размещены здесь.
